\chapter{Tinjauan Pustaka}

% ================================================================
% PANDUAN PENULISAN BAB II — TINJAUAN PUSTAKA
% Bab ini memuat landasan teoretis yang menunjang penelitian Anda.
% Susun subbab sesuai kebutuhan topik, kemudian kembangkan setiap
% poin menjadi narasi akademis Bahasa Indonesia yang didukung sitasi
% (key referensi sudah tersedia di bibliography/references.bib).
% ================================================================

\section{Definisi dan Konsep Dasar}

% Jelaskan definisi istilah/topik utama penelitian, baik menurut
% literatur maupun klasifikasi/standar yang berlaku. Kemukakan
% mengapa topik tersebut penting untuk dikaji.

\section{Klasifikasi}

% Uraikan sistem klasifikasi yang relevan dengan topik (misal
% berdasarkan etiologi, anatomi, waktu awitan, atau derajat keparahan).
% Tekankan implikasi klinis dari masing-masing klasifikasi terhadap
% tata laksana maupun prognosis.

\section{Epidemiologi}

% Sajikan data epidemiologi topik penelitian:
% - Angka prevalensi/insidensi global dan nasional.
% - Distribusi menurut karakteristik demografis (usia, jenis kelamin).
% - Faktor risiko dan mekanisme yang paling sering ditemukan.
% - Pergeseran beban penyakit antar-region, bila relevan.

\section{Etiologi dan Patofisiologi}

% Jelaskan penyebab dan mekanisme terjadinya penyakit/masalah yang
% diteliti. Bila memiliki beberapa mekanisme, uraikan tiap mekanisme
% dalam subbab/subpoin terpisah agar lebih runtut.

\section{Faktor Risiko dan Prediktor}

% Sebutkan faktor-faktor yang meningkatkan risiko atau menjadi
% prediktor luaran klinis, lengkap dengan sitasi pendukung. Sajikan
% dalam narasi atau tabel perbandingan bila jumlahnya banyak.

\section{Manifestasi Klinis dan Diagnosis}

% - Keluhan dan tanda klinis khas.
% - Pemeriksaan fisik dan penunjang yang digunakan untuk menegakkan
%   diagnosis (mis. pemeriksaan khusus, laboratorium, pencitraan).
% - Urutan/strategi diagnosis dan indikasi pemeriksaan penunjang.

\section{Tata Laksana}

% Uraikan prinsip penatalaksanaan berdasarkan bukti terkini, disusun
% sesuai tahapan (konservatif, medikamentosa, tindakan, bedah).
% Jelaskan indikasi, kelebihan, dan keterbatasan tiap pilihan terapi.

\section{Penelitian Terkait}

% - Bahas studi-studi sebelumnya dari berbagai populasi yang relevan
%   dengan topik penelitian Anda.
% - Tunjukkan persamaan, perbedaan, serta kesenjangan yang masih ada.
% - Tekankan keterbatasan data pada populasi lokal sebagai justifikasi
%   penelitian Anda.

\section{Kerangka Teori, Kerangka Konsep, dan Hipotesis}

% - Kerangka teori: skema alur hubungan antarvariabel berdasarkan teori.
% - Kerangka konsep: skema variabel yang diteliti pada penelitian ini
%   (dapat dibuat sebagai gambar/diagram, lihat contoh di Bab IV).
% - Hipotesis: pernyataan dugaan hasil (diisi bila penelitian bersifat
%   inferensial/analitik; untuk penelitian deskriptif bagian ini
%   dapat disesuaikan).

\section{Ringkasan}

% Rangkum poin-poin kunci tinjauan pustaka sebagai jembatan menuju
% Bab III, misalnya: beban masalah yang tinggi, mekanisme yang
% multifaktorial, serta kesenjangan data pada populasi lokal yang
% melatarbelakangi perlunya penelitian dilakukan.
